\documentclass[conference]{IEEEtran}
\IEEEoverridecommandlockouts

\usepackage{amsmath}
\usepackage{amssymb}
\usepackage{array}
\usepackage{booktabs}
\usepackage{cite}
\usepackage{graphicx}
\usepackage{hyperref}
\usepackage{listings}
\usepackage{xcolor}

\lstset{
  basicstyle=\ttfamily\footnotesize,
  breaklines=true,
  frame=single,
  columns=fullflexible
}

\begin{document}

\title{Federated Fraud Detection with Scalable Aggregation, Vietnam-Oriented Data Extension, and Fault-Tolerant Recovery}

\author{
\IEEEauthorblockN{Ngoc Bui}
\IEEEauthorblockA{Federated Learning Fraud Detection Project\\
Email: project author}
}

\maketitle

\begin{abstract}
Financial fraud detection is difficult to centralize because transaction data is
sensitive, institution-specific, and legally constrained. This project builds a
federated learning system for fraud detection in which clients train locally and
only model updates are sent to the server. The baseline dataset is IEEE-CIS
Fraud Detection, and the system is extended with a custom transaction adapter
and a Vietnam-style synthetic payment dataset to prepare for future Vietnamese
banking or wallet data. The proposed system uses a residual multilayer
perceptron, a 328-feature transaction schema, target-aware robust federated
averaging, adaptive training schedules, per-round checkpoint recovery, and a Phi
Accrual failure detector utility. Experimental artifacts show that the 10-client
run reached the deployment target with AUPRC 0.7106, AUROC 0.9431, and F1
0.7025 at round 100. A previous 100-client scalability comparison reached the
target at its best round but regressed by the final round, showing why large
client counts require coverage-aware sampling, robust aggregation, update
stabilization, and more total data rather than only splitting one dataset into
more clients. The system is therefore best described as a research prototype
with working training, inference, checkpoint recovery, large-client simulation,
and a concrete path toward Vietnam-focused deployment.
\end{abstract}

\begin{IEEEkeywords}
federated learning, fraud detection, financial technology, fault tolerance,
scalability, Phi Accrual, Vietnam payments, IEEE-CIS
\end{IEEEkeywords}

\section{Introduction}

Digital payments create large volumes of transactional data, but fraud labels
are rare and the most useful signals are often distributed across separate
organizations. A bank may know account history, a payment gateway may know
merchant behavior, a wallet may know device and channel behavior, and an
e-commerce platform may know order context. Centralizing all of this data is
usually unrealistic because of privacy, regulation, business confidentiality,
and operational risk. Federated learning is therefore a suitable architecture:
each organization keeps its data local, trains a local model, and sends model
updates to a central server for aggregation.

The initial version of this project trained on IEEE-CIS Fraud Detection using
three temporal clients. The project was then expanded after review feedback in
four directions:

\begin{enumerate}
  \item avoid relying on a single dataset by adding a partner/custom transaction
        adapter and a Vietnam-style synthetic dataset path;
  \item compare 3, 10, and 100-client federated behavior;
  \item make the server aggregation algorithm explicit and implement it in code;
  \item improve fault-tolerance discussion with checkpoints, round-level failure
        handling, and a Phi Accrual detector utility.
\end{enumerate}

The project is not presented as a finished banking product. It is a working
research prototype that demonstrates the core distributed-system ideas, the
fraud-modeling pipeline, and the limitations that must be solved before
production use.

\section{Problem Statement}

The technical problem is to train and serve a fraud detection model under four
constraints.

\subsection{Privacy and Data Locality}

Raw transaction data should remain on each client. The server should not require
direct access to private transaction records. This is important for financial
institutions because transaction histories, device identifiers, email domains,
and fraud labels are sensitive.

\subsection{Class Imbalance}

Fraud is rare. In the IEEE-CIS baseline, only about 3.5 percent of transactions
are fraud. Accuracy is therefore misleading: a model that approves every
transaction can appear accurate but is useless. The project uses AUPRC, AUROC,
and F1 as the primary metrics, with AUPRC emphasized because it reflects ranking
quality under class imbalance.

\subsection{Non-IID Client Data}

Client data distributions differ. A temporal split means client 0 may see an
earlier fraud pattern while client 9 sees a later one. A real deployment would
be even more heterogeneous because different banks, wallets, merchants, and
regions have different customer behavior. The aggregation algorithm must
therefore tolerate noisy and uneven local updates.

\subsection{Scalability and Failure}

The system must continue operating when clients fail or when the number of
clients increases. A three-client demonstration is not enough to show a credible
distributed system. The project therefore adds 10-client and 100-client run
paths, checkpoint recovery, and a failure-suspicion utility based on Phi
Accrual.

\section{Related Work}

Federated Averaging (FedAvg) is the standard baseline for federated neural
network training. In FedAvg, selected clients train locally and the server
averages their parameters weighted by sample count \cite{fedavg}. Flower
provides a practical framework for building this workflow and exposes
server-side strategies such as FedAvg \cite{flowerstrategy, flowerdocs}.

Fraud detection is a class-imbalanced ranking task. Davis and Goadrich show why
precision-recall analysis can be more informative than ROC analysis when the
positive class is rare \cite{prroc}. This motivates using AUPRC as the main
model-selection metric.

The model architecture borrows two deep-learning ideas. Residual connections
help train deeper networks by learning corrections to a hidden representation
\cite{resnet}. Layer normalization stabilizes hidden activations without using
batch-level running statistics \cite{layernorm}. This matters in federated
learning because client distributions are non-IID; averaging BatchNorm running
statistics across clients can produce statistics that do not match any real
client.

Fault detection in distributed systems is commonly handled with adaptive
failure detectors rather than fixed timeouts. Phi Accrual converts heartbeat
silence into a suspicion score \cite{phi}. SWIM provides scalable membership
and failure detection using gossip-style probing \cite{swim}. This project
implements Phi Accrual as a tested utility and leaves full live membership
integration as future work.

Robust aggregation is related to Byzantine-resilient learning. Krum,
coordinate median, and trimmed mean reduce the impact of outlier or malicious
updates \cite{krum, byzantine}. This project does not claim full Byzantine
security, but it uses robust blending to reduce large-client instability.

Vietnam payment behavior is relevant because public infrastructure includes
wallet payments, payment gateway flows, QR payments, ATM/POS switching, and
interbank transfer rails. Public documentation from MoMo, VNPAY, and NAPAS
motivates the synthetic Vietnam-style fields, but the project clearly separates
synthetic testing from real Vietnamese fraud evidence \cite{momo, vnpay,
napas}.

\section{Proposed Approaches}

\subsection{Overall Architecture}

The system has five layers:

\begin{enumerate}
  \item data ingestion and feature engineering;
  \item local client training;
  \item server aggregation and checkpointing;
  \item offline and zero-shot evaluation;
  \item FastAPI and React transaction-screening demonstration.
\end{enumerate}

The main implementation modules are listed in Table~\ref{tab:modules}.

\begin{table*}[t]
\caption{Main Implementation Modules}
\label{tab:modules}
\centering
\begin{tabular}{p{0.32\textwidth}p{0.60\textwidth}}
\toprule
\textbf{Module} & \textbf{Responsibility} \\
\midrule
\texttt{dataset/load\_ieee\_cis.py} & Loads IEEE-CIS, merges transaction and identity tables, engineers features, normalizes values, and writes client Parquet files. \\
\texttt{dataset/custom\_transaction\_adapter.py} & Maps partner or synthetic transaction columns into the active feature contract. \\
\texttt{dataset/generate\_vietnam\_synthetic.py} & Generates Vietnam-style payment records for system testing. \\
\texttt{src/data/feature\_registry.py} & Defines the 328-feature schema and label contract. \\
\texttt{src/model/fraud\_mlp.py} & Implements the residual MLP fraud classifier. \\
\texttt{src/client/client.py} & Implements Flower client training, loss functions, sampling, evaluation, and metrics. \\
\texttt{src/server/aggregation.py} & Implements target-aware weighting, robust blending, and update stabilization. \\
\texttt{src/server/strategy.py} & Implements the Flower server strategy, client sampling, checkpoint persistence, and training summaries. \\
\texttt{src/server/failure\_detector.py} & Implements Phi Accrual failure suspicion. \\
\texttt{scripts/run\_local\_flower.py} & Runs local Flower infrastructure, clients, and a submitted run. \\
\texttt{scripts/run\_virtual\_federated.py} & Runs many virtual clients in one process for scalable local experiments. \\
\texttt{api/main.py} and \texttt{app/} & Serve the selected model and provide the transaction-screening GUI. \\
\bottomrule
\end{tabular}
\end{table*}

\subsection{Data Partitioning}

The IEEE-CIS dataset is sorted by transaction time. For \(K\) clients, client
\(k\) receives a contiguous temporal slice:

\begin{equation}
D_k=\{x_i \mid \lfloor kN/K \rfloor \leq i < \lfloor (k+1)N/K \rfloor\}.
\end{equation}

Temporal partitioning is intentionally harder than random splitting because it
reduces future leakage and creates more realistic distribution shift.

\subsection{Feature Engineering}

The active schema is named \texttt{fraud-history-scalable} and contains 328
features. The important feature groups are shown in Table~\ref{tab:features}.

\begin{table}[t]
\caption{Feature Groups}
\label{tab:features}
\centering
\begin{tabular}{p{0.30\columnwidth}p{0.58\columnwidth}}
\toprule
\textbf{Group} & \textbf{Purpose} \\
\midrule
Amount and volume & Capture transaction size, recent spending, and amount pressure. \\
Time features & Encode hour, day, risky-hour flag, and cyclic time behavior. \\
Identity and device & Capture card, email, address, device, and product signals. \\
Frequency encoding & Represent repeated identifiers without exposing raw strings. \\
Historical behavior & Use backward-looking count, mean amount, fraud rate, and time-since-previous features. \\
Risk-shape features & Combine velocity, amount deviation, rare identity, and rule-stack behavior. \\
\bottomrule
\end{tabular}
\end{table}

Numeric features are normalized using training statistics:

\begin{equation}
z_j=\frac{x_j-\mu_j}{\max(\sigma_j,10^{-8})}.
\end{equation}

The normalization parameters are saved in
\texttt{config/normalization\_params.json}. External datasets reuse these
statistics during zero-shot evaluation. Re-fitting the scaler on a new dataset
would change the experiment and hide domain shift.

\subsection{Model Architecture}

The model is a residual multilayer perceptron:

\begin{equation}
h_0=W_0x+b_0,
\end{equation}

\begin{equation}
h_{\ell+1}=h_\ell+f_\ell(\operatorname{LayerNorm}(h_\ell)),
\end{equation}

\begin{equation}
\hat{y}=\sigma(g(h_L)).
\end{equation}

The implementation uses 328 input features, hidden dimension 256, three
residual blocks, and a LayerNorm/ReLU classification head. Hidden dimension 256
is a practical compromise: it is wide enough to represent interactions between
amount, time, identity, and history features, but small enough for repeated
client-side training rounds on a laptop. Three residual blocks give enough
depth for nonlinear feature interactions without making federated updates too
large or unstable. LayerNorm is used instead of BatchNorm because each client
has a different local distribution and fraud rate.

\subsection{Client Training Objective}

Each client trains with a hybrid loss combining focal loss and weighted binary
cross entropy. Focal loss is:

\begin{equation}
\mathcal{L}_{focal}
=-\alpha_t(1-p_t)^\gamma\log(p_t),
\end{equation}

where \(p_t\) is the predicted probability assigned to the true class. Weighted
BCE provides a stable baseline for rare positives:

\begin{equation}
\mathcal{L}_{bce}
=-[w y\log(\hat{y})+(1-y)\log(1-\hat{y})].
\end{equation}

The hybrid loss is:

\begin{equation}
\mathcal{L}=(1-\lambda)\mathcal{L}_{focal}+\lambda\mathcal{L}_{bce}.
\end{equation}

FedProx regularization is optionally added to reduce client drift:

\begin{equation}
\mathcal{L}_{client}=\mathcal{L}+\frac{\mu}{2}\|\theta-\theta_r\|_2^2.
\end{equation}

The client also uses a weighted sampler so that fraud examples appear often
enough to produce a useful gradient while still keeping the training
distribution close to the real distribution.

\subsection{Target-Aware Robust Federated Averaging}

The server aggregation algorithm is explicitly implemented in
\texttt{src/server/aggregation.py}. The target metrics are:

\begin{equation}
T_A=0.70,\quad T_R=0.90,\quad T_F=0.70,
\end{equation}

where \(T_A\), \(T_R\), and \(T_F\) correspond to AUPRC, AUROC, and F1. Client
\(i\) receives a capped target progress score:

\begin{equation}
s_i=0.35\min\left(\frac{AUPRC_i}{T_A},1\right)
+0.20\min\left(\frac{AUROC_i}{T_R},1\right)
+0.45\min\left(\frac{F1_i}{T_F},1\right).
\end{equation}

The server then builds a sample and quality-aware weight:

\begin{equation}
\alpha_i=q_i\sqrt{n_i},\qquad
w_i=\frac{\alpha_i}{\sum_j\alpha_j}.
\end{equation}

The square root reduces domination by a single large client while still
respecting sample size. The weighted FedAvg estimate is:

\begin{equation}
\theta^{weighted}_{r+1}=\sum_i w_i\theta_i.
\end{equation}

For larger client counts, the server blends this weighted average with a robust
coordinate statistic:

\begin{equation}
\theta^{proposed}_{r+1}
=(1-\beta)\theta^{weighted}_{r+1}
+\beta\theta^{robust}_{r+1}.
\end{equation}

Here \(\theta^{robust}_{r+1}\) is a coordinate trimmed mean or coordinate
median. This helps when some clients are noisy, stale, or locally overfit.

\subsection{Server Update Stabilization}

Large federations can produce unstable aggregate jumps. The server therefore
clips and damps the update:

\begin{equation}
\Delta_r=\theta^{proposed}_{r+1}-\theta_r,
\end{equation}

\begin{equation}
\Delta'_r=\Delta_r\min\left(1,
\frac{\rho\|\theta_r\|_2}{\|\Delta_r\|_2+\epsilon}\right),
\end{equation}

\begin{equation}
\theta_{r+1}=\theta_r+\eta_s\Delta'_r.
\end{equation}

For 50 or more configured clients, the default scalable settings are server
learning rate \(\eta_s=0.65\), maximum update ratio \(\rho=0.035\), trim ratio
0.10, and robust blend \(\beta=0.25\).

\subsection{Adaptive Training Schedule}

The schedule controller in \texttt{src/server/training\_schedule.py} observes
recent validation loss, AUPRC slope, F1 slope, and target score. It switches
between warmup, learning, plateau, recovery, and refine phases. When training
regresses, it lowers client learning rate, lowers server learning rate,
increases BCE stability, and increases FedProx strength. When training is
stalled below target, it uses a milder recovery schedule. This prevents the
system from continuing aggressive updates after the model has already started
to degrade.

\subsection{Fault Tolerance and Recovery}

At round \(r\), the server saves:

\begin{equation}
C_r=(r,\theta_r,M_r),
\end{equation}

where \(M_r\) contains metadata such as aggregation weights, participating
clients, schedule state, and metrics. On restart, the server selects the latest
compatible checkpoint:

\begin{equation}
r^*=\max\{r\mid C_r \text{ exists and matches the active schema}\}.
\end{equation}

Flower also allows a round to continue when some clients fail and enough client
results remain. The implemented prototype aggregates successful client results,
records failures, and writes per-round checkpoint files under the configured
run folder.

The Phi Accrual detector utility computes suspicion from heartbeat silence:

\begin{equation}
\phi(t)=-\log_{10}(1-F(t-t_{last})).
\end{equation}

If \(\phi(t)\geq\Phi_{threshold}\), the client is suspected. The current system
implements the detector utility and tests it, but it does not yet run a live
heartbeat membership service inside Flower. This is an honest boundary of the
prototype.

\section{Implementation Details}

\subsection{Runtime Stack}

The system uses Python, PyTorch, Flower, NumPy, Pandas, FastAPI, and React. The
latest local training path uses Flower SuperLink/SuperNode style runtime for
process-based experiments. The local SuperLink defaults to in-memory runtime
state to avoid repeated SQLite lock failures during rapid local experiments.

\subsection{Training Modes}

There are three practical training modes:

\begin{enumerate}
  \item \textbf{3-client mode}: smallest reproducible baseline.
  \item \textbf{10-client mode}: stronger scalability demonstration on one
        machine while still keeping enough samples per client.
  \item \textbf{100-client virtual mode}: keeps 100 client partitions but runs
        them inside one Python process to avoid Windows page-file and DLL
        loading limits from starting 100 PyTorch processes.
\end{enumerate}

The virtual path is important. It preserves federated partitioning and server
aggregation logic, but removes operating-system process memory as the dominant
failure source. This makes it a better local-machine test of algorithmic
scalability.

\subsection{Vietnam-Oriented Dataset Extension}

The Vietnam extension is implemented in two layers. First, a custom adapter maps
external columns to the project schema. Second, a synthetic generator creates
Vietnam-style transactions for integration testing. Synthetic fields include
provider, channel, amount in VND, customer id, merchant id, payer bank, receiver
bank, device id, province, account age, short-window velocity, distance, and
chargeback count.

The role of this dataset is not to claim real Vietnam fraud accuracy. Its role
is to prove that the pipeline can accept Vietnam-shaped transaction data and to
prepare for a real anonymized partner dataset.

\subsection{Inference Application}

The API automatically selects the best compatible checkpoint for a configured
model run. The GUI lets an operator score a transaction, compare reliable and
suspicious examples, and inspect the decision threshold, selected model, USD
amount, and currency conversion status. Currency conversion to USD is handled
in the backend with fallback rates so the model input matches the training
schema.

\section{Evaluation and Experimental Results}

\subsection{Metrics}

The project reports AUPRC, AUROC, and F1. F1 is:

\begin{equation}
F1=\frac{2PR}{P+R}.
\end{equation}

AUPRC approximates the area under the precision-recall curve:

\begin{equation}
AUPRC=\sum_k(R_k-R_{k-1})P_k.
\end{equation}

The deployment target is:

\begin{equation}
AUPRC\geq0.70,\quad AUROC\geq0.90,\quad F1\geq0.70.
\end{equation}

\subsection{Dataset Summary}

The IEEE-CIS baseline contains 590,540 transactions with an approximate fraud
rate of 3.5 percent. A recent 10-client processed split contains 59,054
transactions per client. The client fraud rates range from about 2.02 percent
to 4.32 percent, confirming that temporal client splits are non-IID.

\subsection{Training Results}

Table~\ref{tab:results} reports the available local artifacts at the time of
writing. The 10-client result comes from \texttt{results/10\_clients/}. The
100-client result comes from the scalability comparison artifact and is included
because it demonstrates the key scalability issue: the best round met the target
but later rounds regressed.

\begin{table}[t]
\caption{Available Training Artifacts}
\label{tab:results}
\centering
\begin{tabular}{lrrrr}
\toprule
\textbf{Run} & \textbf{Round} & \textbf{AUPRC} & \textbf{AUROC} & \textbf{F1} \\
\midrule
10 clients latest/best & 100 & 0.7106 & 0.9431 & 0.7025 \\
100 clients best & 44 & 0.7068 & 0.9373 & 0.7248 \\
100 clients latest & 100 & 0.6688 & n/a & 0.6903 \\
Virtual smoke & 1 & 0.4179 & 0.8954 & 0.4887 \\
\bottomrule
\end{tabular}
\end{table}

The 10-client run satisfies the target at round 100. The 100-client comparison
shows a different lesson: the model can reach the target, but the final model
can be worse than the best checkpoint. This supports the added design decisions:
best-checkpoint selection, adaptive recovery, robust aggregation, update
stabilization, and training summaries that distinguish latest and best rounds.

\subsection{Scalability Interpretation}

Scalability is not the same as automatic metric improvement. If the same fixed
dataset is split into more clients, each client receives fewer fraud examples.
This increases variance in local training and validation. Therefore, the
reasonable scalability goal is:

\begin{equation}
\text{more clients} \Rightarrow
\text{stable rounds, broad coverage, bounded regression, recoverable failures}.
\end{equation}

Linear metric improvement requires more total data, such as real data from more
Vietnam banks, wallets, gateways, or merchants. Splitting the same IEEE-CIS data
into 100 clients is useful for stress testing the algorithm, but it does not
create new fraud behavior.

\subsection{Zero-Shot and Cross-Dataset Evaluation}

The codebase includes zero-shot external evaluation for datasets such as PaySim,
credit-card fraud datasets, BAF, and Vietnam synthetic data. The correct method
is:

\begin{enumerate}
  \item inspect the external dataset columns, row count, fraud ratio, and entity
        identifiers;
  \item reconstruct only the compatible features;
  \item reuse the IEEE-CIS scaler without re-fitting;
  \item run inference only, with no gradient updates;
  \item report AUPRC and recall at fixed precision.
\end{enumerate}

Large drops in zero-shot performance should not be hidden. They are evidence of
domain shift and justify the next research step: acquiring real Vietnam partner
data or generating stronger synthetic pretraining data that is clearly labeled
as synthetic.

\subsection{Fault-Tolerance Evaluation}

The implemented recovery path can be evaluated by interrupting training,
restarting from the latest compatible checkpoint, and verifying that round
metadata, model parameters, and training summaries continue under the same run
folder. Unit tests also cover the Phi Accrual failure detector behavior.

\section{Teamwork}

Table~\ref{tab:teamwork} gives a detailed workload split. Names should be
updated before final submission if the project has more than one member.

\begin{table*}[t]
\caption{Team Member Contributions}
\label{tab:teamwork}
\centering
\begin{tabular}{p{0.16\textwidth}p{0.64\textwidth}p{0.10\textwidth}}
\toprule
\textbf{Member} & \textbf{Contribution} & \textbf{Workload} \\
\midrule
Member A & Designed and implemented the federated learning server path, target-aware robust aggregation, update stabilization, checkpoint management, 10/100-client scalability configuration, Flower runtime scripts, fault-tolerance documentation, and IEEE report. & 55\% \\
Member B & Implemented or supported data preprocessing, feature inspection, GUI/API demonstration, transaction examples, experiment execution, result collection, presentation preparation, and validation of the user-facing workflow. & 45\% \\
\bottomrule
\end{tabular}
\end{table*}

\section{Discussion}

The project answers the main distributed-system questions as follows.

\subsection{What happens if the central server crashes?}

The server persists per-round checkpoints. After a crash, it reloads the latest
schema-compatible checkpoint and resumes from that global model. This is
checkpoint-based recovery. The algorithm is the maximum compatible checkpoint
selection:

\begin{equation}
r^*=\max\{r\mid C_r \in \mathcal{C},\operatorname{schema}(C_r)=S\}.
\end{equation}

\subsection{How are fault tolerance and scalability handled?}

Fault tolerance is handled by round-level failure handling, successful-client
aggregation, checkpointing, and restart recovery. Scalability is handled by
coverage-aware sampling, virtual 100-client training, robust aggregation,
adaptive schedules, and run-specific output folders.

\subsection{What if one connection fails?}

If enough clients still respond, the round aggregates the successful results and
logs the failure. A production system would add live membership monitoring using
Phi Accrual or SWIM so the server can distinguish slow clients from failed
clients before scheduling future work.

\subsection{How is unfinished work reassigned?}

In federated learning, the unit of work is a round assignment, not a password
range. If a client fails during round \(r\), its update is excluded and the next
round samples another client. The equivalent reassignment rule is:

\begin{equation}
S_{r+1}=\operatorname{Sample}(\mathcal{A}_r \setminus \mathcal{F}_r),
\end{equation}

where \(\mathcal{A}_r\) is the available client set and \(\mathcal{F}_r\) is
the suspected failed set. No completed global state is lost because the previous
checkpoint remains valid.

\subsection{How does the server verify a result is not fake?}

The current prototype validates shape, schema, and metric consistency, then
uses robust aggregation to reduce outlier influence. It does not yet implement
cryptographic proof or secure aggregation. A production extension should add
signed client updates, secure aggregation, differential privacy if required,
and stronger Byzantine defenses.

\section{Limitations and Future Work}

The largest limitation is dataset realism. IEEE-CIS is a strong public baseline,
but it is not Vietnam-specific. The synthetic Vietnam generator is useful for
pipeline testing, not for final accuracy claims. The next important step is an
anonymized Vietnam partner dataset with transaction time, amount, fraud label,
customer or account id, merchant id, device id, channel, and region.

The second limitation is live fault detection. Phi Accrual is implemented as a
utility, but live heartbeat membership is not yet integrated into the Flower
runtime. The production design should add a membership service that tracks
heartbeats, computes \(\phi(t)\), removes suspected clients from scheduling,
and allows them to rejoin after recovery.

The third limitation is adversarial robustness. Robust blending reduces noisy
updates, but it is not a complete defense against malicious clients. Future work
should evaluate Krum, coordinate median, trimmed mean, signed updates, anomaly
scoring on update vectors, secure aggregation, and differential privacy.

\section{Conclusion}

This project implements a federated fraud detection system with a practical
training pipeline, explicit server aggregation algorithm, transaction-screening
application, Vietnam-oriented data extension, checkpoint recovery, and a path
toward 100-client scalability. The strongest current result is the 10-client
artifact that reaches AUPRC 0.7106, AUROC 0.9431, and F1 0.7025. The 100-client
artifact shows why scalability is a real research problem: the best round can
meet the target while later rounds regress. The final direction is clear:
collect or partner for Vietnam transaction data, integrate live failure
membership, strengthen aggregation security, and evaluate scalability with more
total data rather than only more partitions.

\begin{thebibliography}{00}

\bibitem{fedavg}
H. B. McMahan, E. Moore, D. Ramage, S. Hampson, and B. A. y Arcas,
``Communication-efficient learning of deep networks from decentralized data,''
in \emph{Proc. AISTATS}, 2017.

\bibitem{flowerstrategy}
Flower AI, ``Use a federated learning strategy.'' [Online]. Available:
\url{https://flower.ai/docs/framework/tutorial-series-use-a-federated-learning-strategy-pytorch.html}

\bibitem{flowerdocs}
Flower AI, ``Flower Framework Documentation.'' [Online]. Available:
\url{https://flower.ai/docs/framework/index.html}

\bibitem{prroc}
J. Davis and M. Goadrich, ``The relationship between Precision-Recall and ROC
curves,'' in \emph{Proc. ICML}, 2006, pp. 233--240.

\bibitem{resnet}
K. He, X. Zhang, S. Ren, and J. Sun, ``Deep residual learning for image
recognition,'' in \emph{Proc. CVPR}, 2016.

\bibitem{layernorm}
J. L. Ba, J. R. Kiros, and G. E. Hinton, ``Layer normalization,''
arXiv:1607.06450, 2016.

\bibitem{phi}
N. Hayashibara, X. Defago, R. Yared, and T. Katayama, ``The phi accrual
failure detector,'' in \emph{Proc. SRDS}, 2004.

\bibitem{swim}
A. Das, I. Gupta, and A. Motivala, ``SWIM: Scalable weakly-consistent
infection-style process group membership protocol,'' in \emph{Proc. DSN}, 2002.

\bibitem{krum}
P. Blanchard, E. M. El Mhamdi, R. Guerraoui, and J. Stainer, ``Machine learning
with adversaries: Byzantine tolerant gradient descent,'' in \emph{Proc.
NeurIPS}, 2017.

\bibitem{byzantine}
D. Yin, Y. Chen, R. Kannan, and P. Bartlett, ``Byzantine-robust distributed
learning: Towards optimal statistical rates,'' in \emph{Proc. ICML}, 2018.

\bibitem{targetencoding}
D. Micci-Barreca, ``A preprocessing scheme for high-cardinality categorical
attributes in classification and prediction problems,'' \emph{SIGKDD Explor.
Newsl.}, vol. 3, no. 1, pp. 27--32, 2001.

\bibitem{ieeecis}
IEEE-CIS Fraud Detection, Kaggle competition dataset. [Online]. Available:
\url{https://www.kaggle.com/competitions/ieee-fraud-detection}

\bibitem{momo}
MoMo Developers, ``MoMo Developers documentation.'' [Online]. Available:
\url{https://developers.momo.vn/v3/}

\bibitem{vnpay}
VNPAY Sandbox, ``Payment API documentation.'' [Online]. Available:
\url{https://sandbox.vnpayment.vn/apis/docs/}

\bibitem{napas}
NAPAS, ``National Payment Corporation of Vietnam.'' [Online]. Available:
\url{https://en.napas.com.vn/}

\end{thebibliography}

\end{document}
